\section{Сводные результаты}

Общая сводка по всем сценариям приведена в табл.~\ref{tab:validation-baseline}--\ref{tab:validation-failures} в приложении Б. Для удобства также выделен метод, обеспечивший максимальное социальное благосостояние в каждом сценарии (табл.~\ref{tab:validation-winners}).

\begin{table}[htbp]
	\centering
	\caption{Лучший метод по социальному благосостоянию в каждом сценарии}
	\label{tab:validation-winners}
	\begin{tabular}{lll}
\toprule
Сценарий & Лучший метод по SW & SW \\
\midrule
Сценарий 1: Стабильная нагрузка & QMIX & 3.262 $\pm$ 1.350 \\
Сценарий 2: Пиковая нагрузка & MA-VCG-QMIX & 5.868 $\pm$ 1.097 \\
Сценарий 3: Разнородные ресурсы & Heuristic-LoadAware & 3.547 $\pm$ 0.769 \\
Сценарий 4: Отказы узлов & Heuristic-LoadAware & 3.662 $\pm$ 0.719 \\
\bottomrule
\end{tabular}

\end{table}

Групповое сравнение методов по ключевым метрикам показано на рис.~\ref{fig:validation-overview}. На одном рисунке собраны четыре наиболее важных критерия: социальное благосостояние, доля принятых задач, задержка и коэффициент Джини платежей.

\begin{figure}[htbp]
	\centering
	\includegraphics[width=\textwidth]{method_overview}
	\caption{Сравнение методов по ключевым метрикам}
	\label{fig:validation-overview}
\end{figure}

По результатам сводного анализа можно сделать три принципиальных наблюдения.
\begin{enumerate}
	\item \textbf{Гибридный метод уверенно выигрывает в стабильном режиме и при пиковых нагрузках.} В сценарии~1 \textbf{MA-VCG-QMIX} обеспечивает наибольшее социальное благосостояние ($1.417$), минимальную задержку ($97.4$ мс) и лучшие показатели справедливости среди сравниваемых методов. В сценарии~2 гибридный подход даёт наибольшую долю принятых задач ($41.2\%$) и максимальное значение $\overline{SW} = 2.537$.
	\item \textbf{В сценариях структурной сложности текущая реализация гибридного метода остаётся излишне консервативной.} При разнородных ресурсах и отказах узлов наибольшее социальное благосостояние сохраняет аукционный baseline \textbf{MA-VCG}. Это указывает на то, что RL-компонента прототипа ещё недостаточно агрессивно переобучается в средах с резкой сменой режима.
	\item \textbf{Гибридизация улучшает баланс «эффективность--справедливость», но не устраняет его полностью.} На рис.~\ref{fig:fairness-welfare} видно, что методы не доминируют друг друга одновременно по всем критериям: усиление throughput в сценарии пиковых нагрузок приводит к росту неравномерности платежей, а повышение устойчивости к отказам достигается ценой падения $\overline{SW}$.
\end{enumerate}

\section{Анализ результатов по сценариям}

\subsection{Стабильная нагрузка}

Пошаговая динамика для контрольного сценария приведена на рис.~\ref{fig:baseline-dynamics}. В устойчивом режиме гибридный метод демонстрирует наилучшее сочетание эффективности и устойчивости.

\begin{figure}[htbp]
	\centering
	\includegraphics[width=\textwidth]{baseline_dynamics}
	\caption{Пошаговая динамика метрик в сценарии стабильной нагрузки}
	\label{fig:baseline-dynamics}
\end{figure}

Наиболее важные наблюдения здесь следующие:
\begin{itemize}
	\item \textbf{MA-VCG-QMIX} достигает максимального среднего социального благосостояния ($1.417$ против $1.358$ у \textbf{MA-VCG} и $1.171$ у \textbf{QMIX});
	\item гибридный метод одновременно обеспечивает минимальную задержку ($97.4$ мс) и лучшую равномерность платежей ($Gini = 0.492$);
	\item чистый аукцион сохраняет наибольшую долю принятых задач ($37.4\%$), однако делает это ценой худшей справедливости и более высокой задержки.
\end{itemize}

Следовательно, в стационарной среде интеграция аукциона и QMIX действительно улучшает качество решений по сравнению с использованием только одной из компонент.

\subsection{Пиковая нагрузка}

Динамика при перегрузках представлена на рис.~\ref{fig:load-spikes-dynamics}. Именно этот сценарий лучше всего демонстрирует сильную сторону гибридного подхода.

\begin{figure}[htbp]
	\centering
	\includegraphics[width=\textwidth]{load_spikes_dynamics}
	\caption{Пошаговая динамика метрик в сценарии пиковых нагрузок}
	\label{fig:load-spikes-dynamics}
\end{figure}

В условиях всплесков спроса \textbf{MA-VCG-QMIX} показывает:
\begin{itemize}
	\item максимальное социальное благосостояние ($2.537$);
	\item наибольшую долю принятых задач ($41.2\%$);
	\item наибольшее число завершённых задач ($1.631$ в среднем на шаг).
\end{itemize}

В то же время именно здесь проявляется и ограничение прототипа: коэффициент Джини платежей возрастает до $0.729$. Это означает, что при высокой конкуренции за ресурсы экономический слой ещё не обеспечивает желаемого выравнивания платежей. Тем не менее с точки зрения пропускной способности и эффективности гибридный метод в данном режиме оказывается наиболее сильным.

\subsection{Разнородные ресурсы}

Результаты сценария неоднородных ресурсов показаны на рис.~\ref{fig:heterogeneous-dynamics}. Здесь аукционный baseline остаётся более сильным по $\overline{SW}$, чем обучаемые альтернативы.

\begin{figure}[htbp]
	\centering
	\includegraphics[width=\textwidth]{heterogeneous_dynamics}
	\caption{Пошаговая динамика метрик в сценарии разнородных ресурсов}
	\label{fig:heterogeneous-dynamics}
\end{figure}

Основные результаты:
\begin{itemize}
	\item \textbf{MA-VCG} демонстрирует наибольшее социальное благосостояние ($1.645$) и наибольшую долю принятых задач ($36.2\%$);
	\item \textbf{MA-VCG-QMIX} существенно снижает задержку ($101.8$ мс против $124.3$ мс у \textbf{MA-VCG}) и улучшает обе метрики справедливости;
	\item однако гибридный controller в текущей реализации становится слишком осторожным и теряет часть доступного полезного throughput.
\end{itemize}

С практической точки зрения это означает, что в среде с выраженной гетерогенностью узлов аукционный слой уже содержит достаточно сильный эвристический сигнал, а RL-компонента требует дальнейшей настройки, чтобы не ухудшать суммарную эффективность.

\subsection{Отказы узлов}

Сценарий отказов является наиболее жёстким стресс-тестом и показан на рис.~\ref{fig:failures-dynamics}. Он выявляет устойчивость методов к внезапному уменьшению доступной инфраструктуры.

\begin{figure}[htbp]
	\centering
	\includegraphics[width=\textwidth]{failures_dynamics}
	\caption{Пошаговая динамика метрик в сценарии отказов узлов}
	\label{fig:failures-dynamics}
\end{figure}

Здесь получена смешанная картина:
\begin{itemize}
	\item \textbf{MA-VCG} снова даёт наибольшее социальное благосостояние ($1.617$) и максимальную долю принятых задач ($40.2\%$);
	\item \textbf{QMIX} обеспечивает наименьшую задержку ($84.4$ мс) и лучшую равномерность платежей ($Gini = 0.427$);
	\item \textbf{MA-VCG-QMIX} превосходит \textbf{QMIX} по доле принятых задач ($24.3\%$ против $19.7\%$), но уступает чистому аукциону по эффективности.
\end{itemize}

Таким образом, сценарий отказов показывает текущее ограничение предложенного прототипа: при резком структурном сдвиге среды аукционный механизм остаётся более надёжным, чем обучаемый controller. Однако гибридный метод демонстрирует более сбалансированное поведение, чем QMIX, и не деградирует катастрофически после нарушения топологии.

\section{Кривые обучения}

Эволюция метрик по эпизодам приведена на рис.~\ref{fig:learning-curves}. Этот рисунок особенно важен, поскольку он показывает не только итоговый результат, но и характер адаптации политик.

\begin{figure}[htbp]
	\centering
	\includegraphics[width=\textwidth]{learning_curves}
	\caption{Кривые обучения для QMIX и MA-VCG-QMIX}
	\label{fig:learning-curves}
\end{figure}

По кривым обучения можно выделить следующие тенденции:
\begin{itemize}
	\item в \textbf{сценарии стабильной нагрузки} у \textbf{MA-VCG-QMIX} среднее социальное благосостояние растёт с $1.207$ на первых эпизодах до $1.595$ на последних, а TD-ошибка снижается с $0.417$ до $0.065$;
	\item в \textbf{сценарии пиковых нагрузок} гибридный метод также усиливается по мере обучения: acceptance rate растёт с $36.5\%$ до $43.6\%$, а $\overline{SW}$ увеличивается с $2.188$ до $2.584$;
	\item в \textbf{сценарии разнородных ресурсов} learning-сигнал для гибридного метода стабилизируется, но почти не даёт прироста по $\overline{SW}$, что указывает на недостаточную выразительность текущей политики;
	\item в \textbf{сценарии отказов} наблюдается небольшое улучшение по $\overline{SW}$ у \textbf{MA-VCG-QMIX} ($0.730 \rightarrow 0.850$), однако этого недостаточно, чтобы превзойти аукционный baseline.
\end{itemize}

Следовательно, кривые обучения подтверждают, что гибридная схема действительно обучается и сходится в контролируемых условиях, но её адаптационный потенциал в стрессовых режимах пока ограничен архитектурой прототипа и требует дальнейшего усиления.

\section{Дополнительные артефакты анализа}

Для более детального анализа были построены дополнительные визуальные артефакты.

\subsection{Компромисс между справедливостью и эффективностью}

На рис.~\ref{fig:fairness-welfare} показана диаграмма рассеяния в координатах <<коэффициент Джини платежей --- социальное благосостояние>>.

\begin{figure}[htbp]
	\centering
	\includegraphics[width=0.82\textwidth]{fairness_welfare_scatter}
	\caption{Компромисс между эффективностью и равномерностью платежей}
	\label{fig:fairness-welfare}
\end{figure}

Диаграмма показывает, что текущий прототип не имеет единственной доминирующей стратегии для всех режимов. В стабильной среде и при перегрузках гибридный подход смещает систему в область большей эффективности. При отказах и гетерогенности более сильным остаётся аукционный режим, а по равномерности платежей лучшим часто оказывается QMIX без экономической надстройки.

\subsection{Тепловые карты нагрузки}

Тепловые карты распределения нагрузки для гибридного метода приведены на рис.~\ref{fig:heatmap-baseline}--\ref{fig:heatmap-failures} в приложении C. Они позволяют визуально оценить, насколько равномерно controller распределяет работу между узлами.

Тепловые карты подтверждают сделанный ранее вывод: в нормальных режимах гибридный метод распределяет нагрузку достаточно устойчиво, однако при структурных нарушениях среды поведение становится более фрагментарным и требует усиления механизма переобучения.
